\begin{frame}[label=modern01-setup]{시뮬레이션 도구 확인}\relax
\small
Git·Python 3.10 이상·VS Code·Icarus Verilog를 설치한다. Windows PowerShell에서 한 줄씩 실행한다.\par\terminalbox{git --version\\python --version\\iverilog -V\\vvp -V}
설치 경로를 PATH에 추가했다면 VS Code를 다시 연다. Linux·macOS·WSL은 python 대신 python3를 사용한다.\par
\end{frame}

\begin{frame}{v2.0.0 템플릿을 새 이름으로 clone}\relax
\small
아래는 Windows PowerShell 명령이다. 첫 줄 끝의 백틱(backtick)은 명령을 다음 줄로 이어 준다.\par\terminalbox{git clone --branch v2.0.0 `\\  https://github.com/Glaysia/fpga-lab-template.git lab1\_01\_logic\_gates\\cd lab1\_01\_logic\_gates\\git switch -c main\\code LAB1.code-workspace}
태그로 같은 시작 파일을 받는다. git switch는 태그에서 학생 작업용 main 브랜치를 만든다. 다음 회로는 목적지 폴더 이름을 바꾼다.\par
\end{frame}

\begin{frame}{File → New Window}\relax
\small
VS Code 상단 File → New Window. Ctrl+Shift+N도 같은 동작이다.\par\par\vspace{0.18cm}{\centering\includegraphics[width=\textwidth,height=1.6cm,keepaspectratio,trim=27.75bp 522bp 652.5bp 21bp,clip]{assets/modern-01/01-file-menu.png}\par}\vspace{0.16cm}
새 창의 File → Open Workspace from File...을 누른다.\par\par\vspace{0.18cm}{\centering\includegraphics[width=\textwidth,height=1.6cm,keepaspectratio,trim=27.75bp 439.5bp 652.5bp 105bp,clip]{assets/modern-01/01-file-menu.png}\par}\vspace{0.16cm}

\end{frame}

\begin{frame}{프로젝트 하나의 workspace 열기}\relax
\small
\begin{enumerate}\setlength{\itemsep}{0.22cm}
\item 방금 clone한 lab1\_01\_logic\_gates 폴더를 연다.
\item LAB1.code-workspace를 선택하고 Open. 모든 실습의 workspace 파일명은 같다.
\item Explorer에 PROJECT 하나와 src·sim·constraints·tools 폴더가 보이는지 확인한다.
\item src/design.v·sim/tb\_design.v·constraints/pins.xdc는 아직 주석뿐이다.
\item simulation.json은 실행할 RTL·TB 파일과 TB top을 지정하는 설정이다.
\end{enumerate}

\end{frame}

\begin{frame}{추천 확장 설치}\relax
\small
\begin{enumerate}\setlength{\itemsep}{0.22cm}
\item 왼쪽 Extensions 아이콘을 누른다.
\item 검색창에 @recommended를 입력한다.
\item slang의 Install: Verilog/SystemVerilog 편집·진단.
\item VaporView의 Install: VCD 파형 확인.
\item vscode-pdf의 Install: 코드 옆에서 PDF 교안 열기.
\end{enumerate}
WSL 사용자는 WSL 창에서 필요할 경우 Install in WSL을 누른다. 확장 설치와 Python·Icarus 설치는 별개다.\par
\end{frame}

\begin{frame}{src에 RTL 파일 만들기}\relax
\small
\begin{enumerate}\setlength{\itemsep}{0.22cm}
\item Explorer에서 src 왼쪽 화살표를 눌러 펼친다.
\item design.v 우클릭 → Rename → logic\_gate.v 입력.
\item 파일을 두 번 눌러 열고 뒤의 코드 이미지를 보며 전체 코드를 입력한다.
\item 추가 RTL이 있으면 src 우클릭 → New File로 만든다.
\item 이번 회로의 RTL은 src/logic\_gate.v 하나다.
\item File → Save All로 모든 파일을 저장한다.
\end{enumerate}
설계 모듈명: logic\_gate. 파일명과 module 이름을 구분한다.\par
\end{frame}

\begin{frame}{sim에 테스트벤치 만들기}\relax
\small
\begin{enumerate}\setlength{\itemsep}{0.22cm}
\item sim을 펼치고 tb\_design.v 우클릭 → Rename.
\item 새 파일명: tb\_logic\_gate\_modern.sv
\item 뒤의 TB 코드 이미지에 있는 선언·DUT 연결·입력 자극·예상값 검사를 직접 작성한다.
\item wave.vcd를 만드는 dumpfile·dumpvars와 종료 조건까지 작성한다.
\item TB 모듈명: tb\_logic\_gate\_modern
\item File → Save All. TB를 합성용 RTL 폴더에 넣지 않는다.
\end{enumerate}

\end{frame}

\begin{frame}{simulation.json 전체 내용}\relax
\small
Explorer에서 simulation.json을 열고 아래 7행으로 수정한다.\par\shot{student-simulation-json}sources는 설계 파일, testbench는 검사 파일의 경로다.\par\vspace{0.22cm}simulation\_top은 TB의 module 이름이다. 파일 확장자를 붙이지 않는다.\par\vspace{0.22cm}대괄호·중괄호·큰따옴표·쉼표까지 입력하고 File → Save All.\par
\end{frame}
